\section{它有什么用，下一步是什么？}\label{sec:6}

于是，经过一番堪与赫拉克勒斯之功相比的努力，非分歧几何 Langlands 对应如今已成为定理。但正如我一直试图强调的，这项事业的意义并不只在于、甚至主要不在于验证某个特定陈述为真。围绕 GLC 的研究深化了我们对非阿贝尔对偶的理解——它出现在表示论、算术与量子场论中——并揭示出这些彼此遥远的领域背后一种相当令人惊讶的共同图景。最终导向证明的那一长串论文并不试图走一条测地线式的捷径，而是包围这个陈述、围攻它，直到最后破门而入。这个证明毋宁说是我们理解表示论与调和分析的一个重大里程碑：这是第一次证明了一个如此彻底而一般的非阿贝尔对偶陈述。它表明，作为一个共同体，我们已经理解了非阿贝尔对偶的总体特征，包括其中的关键微妙之处与障碍，并且拥有一个足够的观念框架——以因子化为中心——把这种理解带到新的语境中去。确实，证明中的主要配料——除了 opers 与 Kac–Moody 表示的魔法这一关键例外——基本上都是可以搬运的。

\subsection*{6.0.1\quad 分歧的 GLC}

在走得更远之前，我想简要讨论一下“非分歧”这个修饰词的含义。我们到目前为止讨论中出现的颜色，由光滑射影曲线／紧黎曼面 $C$ 上的（de Rham 或 Betti）局部系统给出。很自然地可以考虑更一般的分歧局部系统，即带亚纯平坦联络的丛，或带洞曲面 $C\setminus\{x_{1},\dots,x_{n}\}$ 的基本群的表示。它们作为只在标记点之外定义的 Hecke 作用的本征值出现——例如在带标记点处附加结构的 $C$ 上丛的模空间上就是如此。Bezrukavnikov~[Bez06, Bez16] 证明了一个里程碑式的定理，把几何 Satake 对应推广到驯顺分歧（丛上的抛物结构／联络上的正则奇点）的情形。该定理尤其给出了 GLC 的驯顺分歧版本的精确表述，而这个版本目前一般情形下仍未解决。

要表述（更不用说证明）GLC 的完整分歧推广，需要对局部 GLC 有更多理解；后者以去心圆盘上的局部系统所给出的 Langlands 参数，为环路群的范畴化表示论提供了一个猜想性的谱分解。关于环绕这一猜想的迅速变化着的图景，见~[Gai25]。

\subsection*{6.0.2\quad 约化群的表示论}

GLC 在表示论中居中心地位，它是与约化群 $G$ 相关的表示的谱理论的原型。这个领域的非凡之处在于：同一个 $G$ 有许多不同的化身，却共享同样的结构特征。自 1990 年代以来，本学科的核心是一个基本的表示论三角，由 Lusztig 的一网猜想联系起来~[Car06]：

\begin{enumerate}
\item 量子群 $U_{q}\mathfrak{g}$；
\item 环路群 $\LG$ 与仿射 Kac–Moody 代数 $\widehat{\mathfrak{g}}$；
\item 正特征的约化群与李代数 $G,\ \mathfrak{g}/\mathbb{F}_{q}$。
\end{enumerate}

来自 GLC 的想法——特别是几何 Satake 对应~[MV07] 及其驯顺分歧的精化（Bezrukavnikov 定理~[Bez06, Bez16]）——已经完全改造了这一领域；很不完整的样本见~[BMR08, BM13, RW22, AR18, Ric24, Wil17, FG09b, RY24]（另见 \S6.0.4）。

\subsection*{6.0.3\quad 低亏格的 GLC}

当我们限制到亏格零与亏格一的曲线时，对应的两侧都变成非常具体而经典的对象，而几何 Langlands 给出的对偶描述提供了深刻的洞见。例如在亏格一时，Langlands 参数是彼此交换的矩阵对，其代数几何长期以来难以捉摸~[LNY24]。在带标记或带洞的亏格零曲线上，Langlands 参数变成亚纯常微分方程，或在 étale 背景中变成一元有理函数域的 Galois 表示。在这一（分歧）背景下，GLC 对一类特殊的刚性局部系统的研究尤为重要~[Yun17]，包括 Kloosterman 指数和的算术与 Bessel 方程~[FG09c, HNY13, XZ22]，以及 Schubert 演算的组合学~[LT24]。一个引人注目的应用是 Yun 对某些有限约化群的反 Galois 问题的解决~[Yun14]。反方向上，GLC 所研究的自守谱问题在低亏格时被等同于著名量子可积系统的求解，例如 Gaudin 自旋链与 Calogero–Moser 粒子系统。于是 GLC 为求解量子可积系统提供了一个强有力的范式：其中彼此交换的 Hamilton 量的联合谱由 Langlands 参数给出~[Fre05]。

\subsection*{6.0.4\quad 其他的 GLC}

GLC 既有退化版本（Dolbeault GLC），也有形变版本（量子 GLC），二者即使在非分歧情形下也仍未解决，尽管如 \S\ref{sec:4.1} 所述，Dolbeault 猜想已在一般情形下由~[DP12] 证明。两者都有一个优美的特征：它们在 $G$ 与 $\Gd$ 之间完全对称；而且二者都自然地出自与 GLC 相同的物理源头（四维 $\mathcal{N}=4$ 超 Yang–Mills 的 S-对偶）。关于受 Donaldson–Thomas 与几何不变量理论启发、在整个 Hitchin 基上研究 Dolbeault 猜想新形式的进展，见~[PT25]。量子 GLC 为仿射 Kac–Moody 代数（de Rham）与量子群（Betti）的（非临界）表示论提供了一个整体的几何语境。它对 $G$ 与 $\Gd$ 的对称处理在若干方面影响了“通常”GLC 的发展，这可以追溯到 Feigin–Frenkel 定理的证明——量子 GLC 正是从中演化而来。例如，GLC 证明中用于“满射性”与“单射性”的 Kac–Moody 周期与 Whittaker 周期都有量子变体，它们处于完全平等的地位，事实上被对偶互换。量子 GLC 中替代几何 Satake 对应的本质局部配料，是所谓基本局部等价~[Gai08, CF21, CDR21, GL19a]，它联系仿射 Kac–Moody 代数的表示与量子群的表示，是回溯到 Drinfeld 与 Kohno 从共形场论的 Knizhnik–Zamolodchikov 方程的单值性构造量子群、以及 Kazhdan–Lusztig 等价~[KL93] 这一系列发展中最精细的一个。

\subsection{算术应用}\label{sec:6.1}

几何 Langlands 纲领最初由 Drinfeld 构想，作为经典 Langlands 对应的一个对应物——在那里人们可以受益于代数几何与层论的强大工具，希望由此对算术中的原初问题获得洞见。在过去十年中，这一希望实现的程度或许超出了合理的预期。下面我聚焦于这一互动最成功、最活跃的两个舞台，即局部域与正特征函数域的背景。这些远非全部。我请读者参阅出色的综述~[Zhu23]，以了解 GLC 对局部与整体 Langlands 纲领、乃至更广泛的算术几何（如 Tate 猜想、Beilinson–Bloch–Kato 猜想）中一系列问题持续涌现的应用中的一些亮点。我也请读者参阅~[Gan]，以了解相对 Langlands 对偶（\S\ref{sec:5.2} 中提及）的综述；在那里几何 Langlands 及其物理解释的想法被应用于 $L$-函数理论以及它们作为自守形式积分（周期）的表示。

\subsubsection{局部 Langlands}\label{sec:6.1.1}

局部 Langlands 对应（先前在 \S\ref{sec:1.4} 中遇到过）为局部域 $F$（如 $\mathbb{Q}_{p}$、$\R$ 与 $\mathbb{F}_{q}((t))$）上约化群的表示提供了一个猜想性的谱理论。换言之，它提议把表示等同于某个“谱”上的向量空间族（重数空间之层）。它并且把参数空间或谱辨识为 Langlands 参数空间——$\Gd(\C)$ 中共轭类的汇集，组织成 $F$ 的某种形式的 Galois 群到 $\Gd(\C)$ 的表示。LLC 是整体 Langlands 纲领的关键伙伴，正如谱定理之于傅里叶变换：整体 Langlands 关涉与 $G$ 和 $F$ 相伴的、$G(\mathbb{F}_{v})$（$v$ 为整体域 $F$ 的一个位）在某个特定表示 $L^{2}([G]_{F})$ 上的谱分解（例如 $L^{2}(\mathrm{SL}_{2}(\Z)\backslash\mathrm{SL}_{2}(\R))$）。

LLC 与 GLC 都采取谱定理的形式，用 Langlands 参数的几何来描述（表示或层的）范畴。这与算术拓扑在局部域与二维流形之间提供的平行十分契合。确实，在 GLC 的影响下，LLC 的研究近来被改造为一个范畴陈述，把局部域上群的表示范畴描述为 Langlands 参数上的拟凝聚层范畴——这是近期 IHES Langlands 纲领会议录第二卷的主题~[CGKS25]。得益于 Fargues–Scholze~[FS21] 在 $p$-进情形与 Scholze~[Sch] 在实情形的工作，这一平行变得直接得多。在这些工作中，局部 Langlands 对应被字面地实现为（非分歧！）几何 Langlands 对应的一个情形，只不过所处背景比域上的曲线奇异得多——分别是 Fargues–Fontaine 曲线与 twistor $\PP^{1}$。GLC 中发展出的许多想法——最显著的是因子化与谱作用——已在此起关键作用，特别是在建立局部 Langlands 的自守到 Galois 方向~[FS21] 以及最近 LLC 中 $\ell$-独立性~[Sch25] 时。这为 GLC 的想法与技术（并间接地，来自物理的结构）影响数论提供了一条直接路径。这是一个极其活跃的领域，新的前沿包括对 $p$-进群的 $p$-进解析表示的 Langlands 对应的新兴几何表述。

\subsubsection{函数域的 Langlands}\label{sec:6.1.2}

几何 Langlands 与经典 Langlands 之间还有一种非常不同（也古老得多）的关系，它来自数域与有限域 $\mathbb{F}_{q}$ 上曲线的函数域之间、以及黎曼面与其代数闭包 $\mathbb{F}_{q}$ 上曲线之间的密切类比（Weil 的罗塞塔石碑）\anote{从算术拓扑的视角看，这对应于把二维流形看作某个纤维化三维流形的纤维。}。在这一类比中，自守形式的向量空间作为自守层范畴的去范畴化出现，其机制是 Frobenius 的迹（Grothendieck 的函数–层对应）。这一机制是 Weil 猜想所给出的、有限域上簇的点计数与其代数闭包上上同调之间联系之直接推广。在这一机制下，自守形式的谱分解（经典 Langlands 对应）是自守层的范畴分解（几何 Langlands）的一个影子。

发展 GLC 的原初梦想，是人们或许可以沿着 Weil 的罗塞塔石碑翻译回去，在从黎曼面过渡到有限域上曲线、并最终也许到数域的过程中带着洞见同行。确实，GLC 的发展在很大程度上是在 $\C$ 上与在 $\mathbb{F}_{q}$ 上平行进行的。正因如此，几何 Langlands 的想法对若干进展至关重要，例如 Ngô 对基本引理的证明~[Nad12, Hal12]——算术 Langlands 纲领的重大突破之一；V.~Lafforgue 建立函数域 Langlands 对应一个方向的工作~[Car, Str17]；Gaitsgory–Lurie 对函数域 Weil 玉河数猜想的证明~[GL19b]；以及 Gaitsgory 对 de Jong 关于有限域上单值性猜想的证明~[Gai07]。

\S\ref{sec:5.1.1} 中讨论的~[AGK+20] 的 étale 版 GLC，在~[AGK+21] 中被证明能通过一种精巧的取 Frobenius 迹的方式，蕴含函数域的非分歧 Langlands 猜想的一个显著强化形式。如今特征零的非分歧 GLC 已被理解，钟摆又摆了回来，函数域的许多重大结构性问题看来已在视野之内。最近~[GR25] 证明了正特征非分歧猜想的一大部分。ICM 报告~[Gai25] 与~[Ras] 勾勒了理解函数域上完全分歧 Langlands 对应的愿景，以及证明正特征自守形式理论中一些最重要的未决问题——Ramanujan 猜想与 Arthur 猜想——的愿景。

\subsection{与物理的关系}\label{sec:6.2}

正如我一直强调的，GLC 的吸引力与力量在很大程度上来自它不只邻近算术，也邻近物理。这里我快速回顾一下。

Kapustin 与 Witten~[KW07]（另见~[Fre10, BZN18]）解释了 GLC 如何自然地作为 Montonen–Olive S-对偶的一个侧面出现——后者是 Maxwell 理论的电磁对称性的一个猜想性非阿贝尔推广。S-对偶是两个四维量子场论之间的非微扰等价，这两个理论分别是两个 Langlands 对偶李群的 Yang–Mills 理论的极大超对称（$\mathcal{N}=4$）扩张。为了得到几何 Langlands，我们从这个极其丰富的陈述中丢弃越来越多的信息。首先我们过渡到理论的一个扭变，即通过过渡到合适微分的上同调（由超对称结构提供），丢掉理论所有对度量的依赖（例如动力学）。（要体会这一过程的猛烈程度，可以看它在 Maxwell 理论（$G=\mathrm{GL}_{1}$ 情形）中的类比：那相当于忘掉电磁学的一切，只留下 Gauss 定律！）有一族依赖于参数 $\Psi\in\PP^{1}$ 的微分，S-对偶给出两个四维 TQFT 之间的等价 $A_{G}\simeq B_{\Gd}$：一个与 $G$ 相伴的“自守”理论（四维 A-模型，$\Psi=0$）和一个与 $\Gd$ 相伴的“谱”理论（四维 B-模型，$\Psi=\infty$）。

然后我们进一步简化，只考虑形如 $C\times\Sigma$ 的四维流形上的理论，其中 $C$ 是固定的黎曼面\anote{构造的不同侧面分别依赖于 $C$ 上的复结构，或仅依赖于其底层的二维流形。}。换言之，我们从一个四维规范场论中提取出一个从黎曼面到二维拓扑量子场论的指派。最后，我们效法同调镜像对称与扩张 TQFT 的做法，对每个这样的二维 TQFT 只聚焦于其边界条件的范畴。于是我们最终得到一个从黎曼面到范畴的指派（这是极其丰富且已被深入研究过的二维共形场论概念的一个高阶类比）。

Kapustin 与 Witten 发现，把这一系列简化施加于物理的 S-对偶猜想，所产生的范畴恰好就是几何 Langlands 对应中出现的那些范畴；而变动 $\Psi$ 则导向量子 GLC（GLC 的 Betti 形式正是作为使这一与 TQFT 的关系显式化的尝试而出现的）。他们的工作以及此后二十年的工作表明：当我们关注越来越多被遗忘的物理结构时，几何 Langlands 纲领的所有关键特征都会浮现，此外还有非凡的新特征。

第一，从四维规范场论（乃至六维量子场论）出发的物理推导，提供了一个统一的框架，用以理解与二维共形场论、经典与量子可积系统、镜像对称、拓扑中的规范场论不变量以及超对称规范场论的 Seiberg–Witten 几何之间的众多联系（其中许多早于~[KW07]）。

第二，我们发现几何 Langlands 只是关于四维理论之陈述的二维部分，因此（猜想性地）推广为涉及最高到四维的各种流形（包括带边界、角与其他缺陷者）的陈述，而这些陈述被 TQFT 的代数结构非常紧密地控制。B-模型 TQFT $B_{\Gd}$ 已得到较好理解，而 A-侧则颇为神秘。然而，得益于 GLC 的证明，我们现在知道（Betti）自守层范畴 $A_{G}(C)=\Shv_{\Nilp}(\Bun_{G}(C))$ 具有非凡的 TQFT 性质：例如，它们带有 $C$ 的映射类群的作用，并满足简单的剪贴规则。最引人注目的是，由 GLC 的证明可以推出，存在一个良定义的向量空间 $A_{G}(M)$，即“三维流形 $M$ 上的自守形式”，它可以用 Heegaard 分解之类的技术、由曲面上的自守层范畴计算出来。它还满足一个 Langlands 对偶，把它等同于特征簇 $\Loc_{\Gd}(M)$（即 $\pi_{1}(M)$ 到 $\Gd$ 的表示空间）上函数空间的（叠化、导出化的）增强。我们直接对这些空间知之甚少，但在形变到量子 GLC 之后，这些空间变成与量子群相伴的三维流形 skein 模——一个被深入研究过的对象，见综述~[Jor23]——从而得到关于 skein 模的 Langlands 对偶预测。

几何 Langlands 的范畴复杂性让位于向量空间之间具体同构的另一个地方，是解析 Langlands 对应~[EFK21]：这是量子可积系统谱分解的一个非凡理论，它在物理上（遵循~[GW24]）作为等价 $A_{G}\simeq B_{\Gd}$ 在一个特殊三维构型上的实例出现。

第三，物理展示了对应两侧之间完全的对称性——毕竟，我们从同一个四维量子场论作用于两个 Langlands 对偶群出发，这一对称性被过程中所作的选择破坏，但可以通过变动那些选择而恢复，然后又被用来增进我们对两侧的理解。这一对称性在数论中完全不明显（那里模形式与 Galois 表示似乎来自不同的宇宙），但它提示了新的视角（例如在相对 Langlands 纲领 \S\ref{sec:5.2} 的框架中）。

最后，超对称量子场论的物理比拓扑场论的物理丰富得多。特别是，它让我们能区分出一类对象——BPS 膜——它们保持比用于定义拓扑理论的微分更多的超对称。这种丰富性在数学上反映为 Hodge 理论的刚性结构。在黎曼面 $C$ 上，这导致希格斯丛与局部系统的模空间从全纯辛几何增强为超 Kähler 几何。在 $\Bun_{G}$ 上，这导致层论从 D-模增强为 Hodge 模或 twistor 模。（这些结构在 \S\ref{sec:4.2} 中遇到的 Donagi–Pantev–Simpson 的工作中起关键作用。）这提示了与 motive 理论的联系。Motivic 层像这些超对称膜一样，有许多不同的实现（de Rham、Betti、Dolbeault、étale）。它们近来在 GLC 与 LLC 中崭露头角~[Sch25]。更广泛地说，理解 motive 本身（作为 Galois 表示的增强）是整个 Langlands 纲领的主要动机之一。人们不禁梦想，超对称与 S-对偶的物理或许能对这些深刻问题提供某种洞见。
